\section{Further exact research questions}
\label{research:section}

The results above change the starting point of the next investigation.
The all-zero-block generator and the single marked transverse derivative
need no further conjectural existence statement. The unresolved issues
concern additional factors, several marked positions, arithmetic
reduction, and analytic continuation beyond the domains proved here.
The questions below specify an object to compute and a criterion for
progress. They are proposals, not assertions that the suggested
reductions hold.

\subsection{Higher products and logarithmic moments}

\paragraph{Q1. Four and more even-order tails.}
For an even number $2d\geq4$ of even orders $a_j\geq2$, consider
\[
 F_{\boldsymbol a}(s)=\sum_{n\geq0}(n+1/2)^{-s}
                         \prod_{j=1}^{2d}\zeta(a_j,n+1).
\]
The centered expansion has exponents
$2d-\sum_j a_j-2k$, so the same subtraction argument makes all even
spectral integers regular. The open arithmetic problem is to evaluate
every $F_{\boldsymbol a}(-2M)$ in an explicit finite list of
multiple-zeta coordinates whose size is independent of $M$.
The first useful case is $(a_1,a_2,a_3,a_4)=(2,2,2,4)$.
Finite summation by parts and ordering the surviving indices give a
concrete route, but the two-factor Euler reduction does not by itself
remove every higher-depth coordinate. A satisfactory theorem should
give the finite coordinate map first and prove every subsequent
reduction. For an odd number of even tails, possible poles occur at
even spectral integers; that problem additionally requires a specified
Laurent finite part and is a different extension.

\paragraph{Q2. Which logarithmic remainders cancel?}
Corollary~\ref{mt:spectralderivatives} gives a convergent exact formula
for every $F_{a,b}^{(r)}(-2M)$, including its finite Hurwitz correction.
The remaining question is whether there are nontrivial, explicitly
specified rational linear combinations of these derivatives whose
logarithmic remainder sums cancel or reduce to endpoint data.
Begin with $r=1$, $(a,b)=(2,4)$ and $(2,2)$, and retain the
$-\log2/6$ term in \eqref{mt:derivative-example}.
One should seek an identity at the level of the summands or a
parameter-dependent kernel; an integer relation among a few decimal
values would only suggest a target. Independent differentiation of a
tail order, such as $\partial_a F_{a,b}$ before specializing $a$,
is another problem because it differentiates both the Hurwitz tail and
its asymptotic coefficients.

\Needspace{5\baselineskip}
\paragraph{Q3. A smallest proved span for each pair of orders.}
Corollary~\ref{mt:closure} supplies at most $\max(a,b)$ coordinates for
the entire sequence. Determine all reductions of that list obtainable
from stated algebraic identities, and give an algorithm producing a
smaller \emph{proved} spanning list. There are two distinct questions:
the rank of explicit rational coefficient vectors in a chosen formal
zeta-product alphabet, and arithmetic independence of the resulting
numbers. Only the first is a finite linear-algebra computation without
additional period theorems. A uniform description for $(2,2q)$ would
be a useful initial result.

\subsection{Several moving positions and primitive calculus}

\paragraph{Q4. Two independently marked positions.}
Construct a jointly continued gap formula for
\[
 \left.\partial_u^r\partial_v^q
       \HH_{a,b}(s,u,v,t)\right|_{u=-m,\,v=-n},
 \qquad m,n\geq0,\quad r,q\geq1,
\]
and identify a finite family of decorated lower-depth objects that
remains sufficient after arbitrary unmarked nonpositive deletions.
The adjacent case $m=n=0$, $r=q=1$ has a precise finite-gap guide.
For two surviving positive real lattice values $y>z$ in a finite
unit-lattice interval, put
\[
 P_r(y,z)=\ell_{0,r}(z+1)-\ell_{0,r}(y).
\]
The elementary symmetric-polynomial identity gives
\[
 \sum_{y>i>j>z}\log i\,\log j
                 =\frac12\bigl(P_1(y,z)^2-P_2(y,z)\bigr).
\]
Here $P_1$ is the negative sum of the gap logarithms and $P_2$ is
their squared-logarithm sum. This finite identity suggests the shape
of a Gamma--Stieltjes formula. A theorem at arbitrary endpoints must
still identify its decorated continuation with the prescribed
relative quotient. The separate-position and unequal-center cases
may require more than products of single-gap primitives.

\paragraph{Q5. Ordered versus symmetrized transverse data.}
Equation~\eqref{tr:insertionderivative} evaluates the sum of the three
possible insertion positions using depth-two data. Determine an exact
decomposition of several marked derivatives into their symmetrized
part and the remaining ordered part. At two marked positions,
stuffle products provide explicit relations but do not identify all
ordered differences. A useful result would describe a finite basis
of these ordered differences modulo the stated stuffle and deletion
relations, with the endpoint normalization retained throughout.

\paragraph{Q6. Endpoint differentiation and normalized antiderivatives.}
Combine the transverse calculus with differentiation in $a$ and $b$.
The identities
\[
 \partial_x\zeta(u,x)=-u\zeta(u+1,x),\qquad
 \mathcal B_m'(x)=\mathcal B_{m-1}(x)
\]
suggest a closed transport system connecting moving harmonic orders,
polygamma factors, and the balanced primitive ladder. Determine that
system for a depth-three word with one marked position, then integrate
it between specified endpoint values. The target is a formula whose
Gamma or Barnes constants are fixed by the same relative quotient,
rather than by separately chosen indefinite integrals. At $r\geq2$,
the analogous system involves generalized Stieltjes primitives and
should state their additive normalization explicitly.

\subsection{Rational shifts and continuation of the generators}

\paragraph{Q7. Common rational shifts and colored tails.}
For rational $c>1/2$, replace the moment family by
\[
 \sum_{n\geq0}(n+c-1/2)^{-s}
                 \zeta(a,n+c)\zeta(b,n+c).
\]
Its centered parity remains available, but the removed polynomial
terms now have continued values $\zeta(-2j,c-1/2)$, which generally
do not vanish. Establish the exact analogue of the finite Bernoulli
formula with all these counterterms. Then determine the appropriate
colored or cyclotomic multiple-zeta reduction for rational $c$.
Unequal shifts in the two tails should be treated separately, since
they no longer share the same centered expansion. Any differentiation
of rational multiplication formulas must keep the conductor powers
and their logarithmic derivatives.

\paragraph{Q8. Rational endpoint Gamma--Stieltjes identities.}
Specialize the finite endpoint formula
\eqref{tr:negativeformula} to $a,b\in\Q_{>0}$ and apply Hurwitz
distribution and character projections. For example,
\[
 \sum_{j=1}^{q}\zeta(s,j/q)=q^s\zeta(s)
\]
gives a completely specified hierarchy after differentiation at
$s=-m$. Determine which linear combinations of transverse weighted
tails reduce to derivatives of Dirichlet $L$-functions and logarithms
of Gamma or Barnes values, and which still require ordered coordinates.
The equality of the finite endpoint expressions provides an exact
starting point; no conjectural independence of Stieltjes constants is
needed for this transport step.

\paragraph{Q9. Singular hyperplanes beyond the common twist tube.}
Theorem~\ref{zbg:entire} proves joint holomorphy on $\Omega_k$.
The kernel identifies the only remaining candidate singular
hyperplanes as nonzero exponential periods of consecutive sums.
Determine their residues or local branching, first at depth two,
including further cancellations at specific nonzero periods.
For integral $a-b$, the function $z^{a-b}$ is single valued near
every nonzero node, so all finite node collisions are removable;
for $a-b=N\geq0$ the finite exponential sum gives the same conclusion.
For nonintegral endpoint differences,
a precise local continuation theorem would explain which parts of
the unit decoration polydisc can be enlarged and on which sheets.
It should treat the complete relative expression, since separately
continued Lerch summands can conceal cancellations.

\paragraph{Q10. Functional identities for the all-moment generator.}
The ODE and polylogarithmic representation in
Theorem~\ref{mt:generator} reduce each fixed pair of orders to
finitely many anchored primitives. Seek recurrences linking different
pairs $(a,b)$, reflection or continuation identities in $t$, and exact
special values at nonzero arguments where the polylogarithms have a
useful algebraic argument. The common radius $2\pi$ also makes the
nearest singularities at $\pm2\pi i$ accessible. Their local
expansions may give exact coefficient relations or full asymptotic
expansions of the moment sequence. Estimates alone are a secondary
goal here; the primary target is a new identity between the generators.

\Needspace{10\baselineskip}
\subsection{Integration with the existing conjectural program}

\paragraph{Q11. An analytic bridge to the Gaussian $S_6$ and $S_8$ candidates.}
The source manuscript distinguishes its formal Cayley or period
coordinates from proved evaluations. The new generators should be
compared with the existing Gaussian kernels only through an explicit
analytic transformation: a convergent integral identity, a justified
parameter deformation, or a proved functional equation. Determine
whether such a transformation sends a specific conjectural remainder
to a Gamma--Stieltjes gap or a mixed-tail coordinate evaluated here.
Until that bridge is proved, the present exact results supply neither
a proof nor a refutation of the $S_6$ or revised $S_8$ formulas.

\paragraph{Q12. Formalization of the exact algebra before the analytic layer.}
The divided-difference product identity, finite endpoint composition,
strict-gap binomial coefficients, Bernoulli deletion, and finite
summation-by-parts formulas are suitable first formalization targets.
Their statements should use exact rational or symbolic arithmetic and
retain the boundary terms. The analytic layer then needs explicitly
stated normalized-Mellin continuation, local normal convergence,
and Hurwitz asymptotics. The accompanying programs are reproducible
checks of finite instances; they do not replace those quantified
theorems. A useful integration milestone is a mechanically checked
finite algebraic core with the analytic lemmas listed as precise
dependencies.

\medskip
For an efficient next investigation, the most concrete starting points
are the four-tail example in Q1, the adjacent two-marked gap in Q4,
and the rational-shift extension in Q7. Each starts from a finite
identity or a convergent kernel already available here. They also test
three different potential obstructions: higher arithmetic depth,
nonlinear ordered primitives, and nonvanishing endpoint counterterms.
